% Standalone theorem insert for W33 papers.
% File: paper/w33_signed_nullspace_point_line_relations_insert.tex

\section{The Signed Point--Line Nullspace Theorem}
\label{sec:signed-point-line-nullspace}

Let $A\in\{-1,0,+1\}^{160\times1620}$ be the signed flag-quadrangle phase matrix.
Rows are point-line flags of $W(3,3)$ and columns are ordinary quadrangles.  We have
already seen that
\[
        \operatorname{rank}(A)=81.
\]
Since there are $160$ flag rows, the left nullity is
\[
        160-81=79.
\]

\begin{theorem}[Signed Point--Line Nullspace Theorem]
For each point $p$ of $W(3,3)$, the four flags through $p$ satisfy a signed four-term
relation among the rows of $A$.  For each totally isotropic line $L$, the four flags
on $L$ satisfy a signed four-term relation among the rows of $A$.  Thus there are
\[
        40+40=80
\]
local signed point/line relations.  These relations have one global dependence, so
that their span has dimension $79$.  Moreover,
\[
        \operatorname{span}\{\text{signed point/line four-flag relations}\}
        =\ker(A^T)=\ker(AA^T).
\]
\end{theorem}

\begin{proof}[Finite certificate]
The certificate constructs $A$ from the signed $\mathbb F_3$ flag-quadrangle pairing.
For every point $p$, it searches the $2^4$ possible sign assignments on the four flags
through $p$ and finds exactly two, differing only by global sign, that annihilate all
columns of $A$.  The same search is performed for every line $L$.  The resulting
$160\times80$ relation matrix has rank $79$, and direct multiplication gives
\[
        A^T R=0,
        \qquad
        AA^T R=0.
\]
Since $\dim\ker(A^T)=160-81=79$, these local relations span the full left nullspace.
\end{proof}

\begin{corollary}[Explicit projector decomposition]
The signed phase projector decomposes the $160$-dimensional flag surface as
\[
        \mathbb R^{160}
        =\operatorname{im}(AA^T)\oplus\ker(AA^T),
\]
with
\[
        \dim\operatorname{im}(AA^T)=81,
        \qquad
        \dim\ker(AA^T)=79.
\]
The kernel is precisely the local signed point/line relation space.  Therefore the
phase matrix kills local gauge relations and leaves the protected $81$-dimensional
homology image.
\end{corollary}
